\documentclass[Total.tex]{subfiles}
\begin{document}
	Two of theories should be contained here:
	\begin{itemize}
		\item Measure Theory;
		\item Real Analysis;
	\end{itemize}
	\chapter{Measure}
		\section{Construction and Extensions of Measures}
			\subsection{Algebras, Semi-Algebras}
				\subsubsection{Algebras}
					\begin{Def}
						Let $X$ be arbitrary set, then a collection $\mathcal{A}$ of subsets of $X$ is an algebra on $X$ if:
						\begin{itemize}
							\item $X,\emptyset\in \mathcal{A}$;(A I)
							\item \begin{gather*}
								A,B\in \mathcal{A}\Rightarrow \begin{cases}
									A\cap B\in \mathcal{A}\qquad (\mbox{A II})\\
									A\cup B\in \mathcal{A}\qquad(\mbox{A III})\\
									A-B\in \mathcal{A}\qquad (\mbox{A IV})
								\end{cases}
							\end{gather*}
						\end{itemize}
					\end{Def}
					
					\begin{proposition}
						For algebra axiom, define
						\begin{gather*}
						 	B\in \mathcal{A}\Rightarrow X-B\in \mathcal{A}\qquad\mbox{(A iv)}
						\end{gather*}
						
						Then we have
						\begin{gather*}
							\{(AI),(AII),(AIII),(AIV)\}\quad\mbox{equivalent to }\{(AI),(AII),(AII),(Aiv)\}
						\end{gather*}
					\end{proposition}
					\begin{proof}
						内容...
					\end{proof}
					
					\begin{proposition}
						TFAE:
						\begin{itemize}
							\item $X\in \mathcal{A}$;
							\item There exists a set $E\in \mathcal{A}$ which is called the \textbf{unit} of the algebra, if
							\begin{gather*}
								A\cap E=A\qquad \forall A\in \mathcal{A}
							\end{gather*}
						\end{itemize}
					\end{proposition}
					\begin{proof}
						$\Rightarrow$ Choose $E=X$;
						
						$\Leftarrow$ 
					\end{proof}
					
					And $\emptyset\in \mathcal{A}$ is not necessary. Notice 
					\begin{gather*}
						\emptyset=X-X\in \mathcal{A}
					\end{gather*}
					
					\textbf{Remark} Such algebra is actually a kind of boolean algebra. We have following correspondence
					\begin{gather*}
						A\cup B\qquad a\vee b\\
						A\cap B\qquad a\wedge b\\
						X-A\qquad \neg a\\
						\emptyset\qquad 0\\
						X\qquad 1\\
						A\subseteq B\qquad a\leq b
					\end{gather*}
				\subsubsection{$\sigma$-Algebra}
					\begin{Def}
						An algebra of sets $\mathcal{A}$ is called a \textbf{$\sigma$-algebra} if for any sequence of sets $(A_n)_{n\in \omega}$, we have
						\begin{gather*}
							A_n\in \mathcal{A}\Rightarrow \cup A_n\in \mathcal{A}\qquad \mbox{(SA I)}
						\end{gather*}
					\end{Def}
					Then
					\begin{Def}
						A pair $(X,\mathcal{A})$ consisting of a set $X$ and a $\sigma$-algebra $\mathcal{A}$ is called a \textbf{measurable space}.
					\end{Def}
					
					\begin{Def}
						Let $X$ be a set. For any family $\mathscr{F}$ of subsets of $X$, there exists a unique $\sigma$-algebra generated by $\mathscr{F}$, denoted by symbol
						\begin{gather*}
							\sigma(\mathscr{F})
						\end{gather*}
						
						which is the smallest $\sigma$-algebra generated by $\mathscr{F}$. 
					\end{Def}
					\begin{proposition}
						Let $X$ be a set. 
						
						Then, for any family of subsets of $X$, there exists a unique $\sigma$-algebra generated by $\mathscr{F}$.
						
						In addition, there exists a unique algebra generated by $\mathscr{F}$. 
					\end{proposition}
					\begin{proof}
						Such algebra is just the intersection of those $\sigma$ algebras, which contains $\mathscr{F}$. 
					\end{proof}
					
				\subsubsection{Morphisms}
					
					\begin{proposition}
						Let $(A_\alpha)$ be a family of subsets of a set $X$, and letr $f:E\rightarrow X$ be an arbitrary mapping of a set $E$ to $X$. Then,
						\begin{gather*}
							X-\cup_{\alpha\in \Lambda} A_\alpha=\cap\qquad X-\cap_{\alpha\in  \Lambda} A_\alpha=\cup_{\alpha\in \Lambda} (X-A_\alpha)\\
							f^{-1}(\cup_{\alpha\in \Lambda} A_\alpha)=\cup_{\alpha\in \Lambda} f^{-1}(A_\alpha)\qquad f^{-1}(\cap_{\alpha\in \Lambda} A_\alpha)=\cap_{\alpha\in \Lambda} f^{-1}(A_\alpha)
						\end{gather*}
					\end{proposition}
					\begin{proof}
						内容...
					\end{proof}
					
					
					\begin{proposition}
						Let $E,X$ be sets, with function
						\begin{gather*}
							f:E\rightarrow X
						\end{gather*}
						\begin{itemize}
							\item 
							Let $\mathcal{A}$ be a $\sigma$-algebra of subsets of a set $X$, 
							
							Then, the class $f^{-1}(\mathcal{A})$ of all sets of the form $f^{-1}(A),A\in \mathcal{A}$, is a $\sigma$-algebra in $E$.
							\item Let $\mathcal{B}$ be a $\sigma$-algebra of subsets of a set $X$. 
							
							Then, the class $\{A\subseteq X|f^{-1}(A)\in \mathcal{B}\}$ is a $\sigma$-algebra;
							\item  Furthermore, for a class of subsets $\mathcal{F}\subseteq P(X)$
							\begin{gather*}
								\sigma(f^{-1}\mathscr{F})=f^{-1}(\sigma(\mathscr{F}))
							\end{gather*}
						\end{itemize}
					\end{proposition}
					\begin{proof}
						\begin{itemize}
							\item 
							\item 
							\item 
						\end{itemize}
					\end{proof}
					
					This is actually the measure induced by function $f$. Could be seen as a pullback. 
					
					\begin{Def}
						Take the following notion of induced algebras:
						\begin{gather*}
							f^{-1}\mathcal{A}\qquad f\mathcal{B}
						\end{gather*}
					\end{Def}
					
				\subsubsection{Borel $\sigma$-Algebra}
					\begin{Def}
						The Borel $\sigma$-algebra of $\RR^n$ is the $\sigma$-algebra
						\begin{gather*}
							B(\RR^n)=\sigma(Open(\RR^n))
						\end{gather*}
						
						And take the following notion: $E\subseteq \RR^n$. Then, 
						\begin{gather*}
							B(E)=\{E\cap B|B\in B(\RR^n)\}
						\end{gather*} 
					\end{Def}
					
					\begin{proposition}
						The Borel $\sigma$-algebra of the real line is generated by any of the following classes of sets:
						\begin{itemize}
							\item The collection of all intervals;
							\item The collection of all intervals with rational endpoints;
							\item The collection of all rays of the form $(-\infty,c)$, where $c$ is rational;
							\item The collection of all rays of the form $(-\infty,c]$, where $c$ is rational;
							\item The collection of rays of the form $(c,\infty)$, where $c$ is rational;
							\item The collection of rays of the form $[c,\infty)$, where $c$ is rational;
						\end{itemize}
						
						Finally, the same is true if in place of rational numbers, one take points in \underline{anywhere dense set}.
					\end{proposition}
					\begin{proof}
						Omitted. 
					\end{proof}
					
					\begin{proposition}
						Consider the collection $\mathscr{F}=\{\{x\}|x\in X\}$, then
						\begin{gather*}
							\sigma(\mathscr{F})=\{A\subseteq X|A\mbox{ is at most countable}\}\cup \{A\subseteq X|X-A\mbox{ is at most countable}\}
						\end{gather*}
						
						Thus, $B(\RR^n)<\sigma(\{\{x\}|x\in \RR^n\})$.
					\end{proposition}
					\begin{proof}
						Easy to verify.
					\end{proof}
				\subsubsection{Ring}
					\begin{Def}
						\begin{itemize}
							\item Let $X$ be a set, with $S\subseteq P(X)$. Then, $S$ is a \textbf{semiring of sets} if
							\begin{gather*}
								\emptyset\in S\qquad A\cap B\in S,\forall A,B\in S\qquad 
								\sqcup_{i=1}^n C_i\in S,C_i\in S
							\end{gather*}
							\item A family $R$ of subsets of a set $X$ is called a \textbf{ring} if
							\begin{gather*}
								\emptyset\in R\qquad A,B\in R\Rightarrow A\cap B,A\cup B,A-B\in R
							\end{gather*}
							\item A family $S$ of subsets of $X$ is called \textbf{semialgebra} if
							\begin{gather*}
								\emptyset\in S\qquad A,B\in S\Rightarrow A\cap B\in S\qquad A,B\in S,A\subseteq B\Rightarrow B-A=\sqcup A_i,
							\end{gather*}
							\item A ring is called a \textbf{$\sigma$-ring} if it is closed with respect to countable unions. i.e. 
							\begin{gather*}
								\cap A_n\in \mathcal{S}
							\end{gather*}
							
							A ring is called a \textbf{$\delta$-ring} if it is closed with respect to countable intersections. i.e.
							\begin{gather*}
								\cup A_n\in \mathcal{S}
							\end{gather*}
						\end{itemize}
					\end{Def}
					
					\begin{proposition}
						\begin{itemize}
							\item Ring$\Rightarrow$ Semiring;
							\item Algebra$\Rightarrow$ Semialgebra;
						\end{itemize}
					\end{proposition}
					\begin{proof}
						\begin{itemize}
							\item 
							\item 
						\end{itemize}
					\end{proof}
					
					\begin{proposition}
						For any semising $S$, 
						the collection
						\begin{gather*}
							R=\{\cup_1^n X_i|X_i\in S\}
						\end{gather*}
						
						forms a ring $R$.
					\end{proposition}
					\begin{proof}
					\end{proof}
					
					And
					\begin{proposition}
						For an algebra $\mathcal{A}$, TFAE:
						\begin{itemize}
							\item $\mathcal{A}$ is $\sigma$;
							\item $\mathcal{A}$ is $\delta$;
						\end{itemize}
					\end{proposition}
					\begin{proof}
						\begin{itemize}
							\item 
							\item 
						\end{itemize}
					\end{proof}
					
					\begin{Def}
						Suppose $\mathscr{B}$ be a $\sigma$-algebra in space $X$, and any set $A\subseteq X$, the class 
						\begin{gather*}
							\mathscr{B}_A=\{B\cap A|B\in \mathscr{B}\}
						\end{gather*}
						
						The $\sigma$-algebra is called the \textbf{trace $\sigma$-algebra}.
					\end{Def}
			\subsection{Measure}
			\subsubsection{Additivity, Subadditivitty;Countable Additivity}
			\begin{Def}Let $A$ be a set, take the notion of power set $P(A)$, and $A_i$ be a family of disjointed sets in $A$. 
				
				Then, a \textbf{real-valued set function} $\mu:\mathcal{A}\rightarrow (-\infty,\infty)$ is defined on a class of sets $\mathcal{A}$ is called
				\begin{itemize}
					\item 
					 \textbf{additive} if
					\begin{gather*}
						\mu(\cup_{i=1}^n A_i)=\sum \mu(A_i)
					\end{gather*}
					
					\item \textbf{subadditive} if
					\begin{gather*}
						\mu(\cup_{i=1}^n A_i)\leq \sum_{i=1}^n \mu(A_i)
					\end{gather*}
					
					\item \textbf{countably additive} if 
					\begin{gather*}
						\mu(\cup_{n=1}^\infty A_n)=\sum_{n=1}^\infty \mu(A_i)
					\end{gather*}
				\end{itemize} 
			\end{Def}
			
			\begin{proposition}
				Let $\mu$ be an additive real set function on an algebra/ ring of sets $\mathcal{A}$, then the following conditions are equivalent:
				\begin{itemize}
					\item The funciton $\mu$ is countably additive;
					\item The function $\mu$ is continuous at zero in following sense: For the following tower
					\begin{gather*}
						A_n\in \mathcal{A}\qquad A_{n+1}\subseteq A_n
					\end{gather*}
					
					Then
					\begin{gather*}
						\mu(\cap_{n=1}^\infty A_n)=0
					\end{gather*}
					
					\item The function $\mu$ is continuous from below. i.e. If $A_n\in \mathcal{A}$ are such that
					\begin{gather*}
						A_n\subseteq A_{n+1},\forall n\in \NN\qquad \cup_{n=1}^\infty A_n\in\mathcal{A}
					\end{gather*}
					
					Then, 
					\begin{gather*}
						\mu(\cup_{n=1}^\infty A_n)=\lim_{n\rightarrow \infty} \mu(A_n)
					\end{gather*}
				\end{itemize}
			\end{proposition}
			\begin{proof}
				\begin{itemize}
					\item 
					\item 
					\item 
				\end{itemize}
			\end{proof}
			
			Furthermore,
			\begin{Def}Let $X$ be a set,
				
				A function
				\begin{gather*}
					\mu_0:\mathcal{A}\rightarrow [0,+\infty]
				\end{gather*}
				\begin{itemize}
					\item If  $\mathcal{A}\subseteq P(X)$ be an ring of sets. $\mu_0$ is called \textbf{premeasure} if
					\begin{itemize}
						\item $\mu_0(\emptyset)=0$;
						\item For a disjointed sequence $(A_n)\subseteq \mathcal{A}$,
						\begin{gather*}
							\cup_{n=1}^\infty A_n\in \mathcal{A}\Rightarrow \mu_0(\sqcup_{n=1}^\infty
							A_n)=\sum_{n=1}^\infty \mu_0(A_n)
						\end{gather*}
					\end{itemize}
					
					\item If  $\mathcal{A}\subseteq P(X)$ is aj algebra of sets.$\mu_0$ is called \textbf{measure} if it is countably additive. 
				\end{itemize}
			\end{Def}
			
			\begin{proposition}
				Let $\mu$ be a nonnegative additive set function on an algebra or ring $\mathcal{A}$:
				\begin{itemize}
					\item If $A,B\in \mathcal{A},A\subseteq B$, then $\mu(A)\leq \mu(B)$;
					\item (Semiadditivity)For a collection 
					\begin{gather*}
						\mu(\cup_{i=1}^n A_i)\leq \sum_{i=1}^n \mu(A_i)
					\end{gather*}
					
					\item (Countably-Semiadditivity) For a collection $\{A_n\}_{n\in\infty}\subseteq\mathcal{A}$, with $\cup_{n=1}^\infty A_n\in \mathcal{A}$, one has
					\begin{gather*}
						\mu(\cup_{n=1}^\infty A_n)\leq \sum_{n=1}^\infty \mu(A_n)
					\end{gather*}
				\end{itemize}
			\end{proposition}
			\begin{proof}
				\begin{itemize}
					\item 
					\item 
					\item 
				\end{itemize}
			\end{proof}
		
			\begin{proposition}
				Let $\mu$ be a nonnegative additive set function on an algebra or a ring $\mathcal{A}$:
				
				\begin{itemize}
					\item If $A,B\in \mathcal{A},A\subseteq B,\mu(A)\leq \mu(B)$;
					\item For any collection $A_1,\dots,A_n\in \mathcal{A}$, one has
					\begin{gather*}
					\mu(\cup_{i=1}^n A_i)\leq \sum_{i=1}^n \mu(A_i)
					\end{gather*}
					\item The function $\mu$ is countably additive precisely, when in addition to the additivity, it is countably subadditive in the following sense: For any sequence $\{A_n\}\subseteq \mathcal{A}$, with $\cup_{n=1}^\infty A_n\in\mathcal{A}$, one has
					\begin{gather*}
						\mu(\cup_{i=1}^\infty A_i)\leq \sum_{i=1}^\infty \mu(A_i)
					\end{gather*}
				\end{itemize}
			\end{proposition}
			\begin{proof}
				\begin{itemize}
					\item Easy to check;
					\item Obtained from (1)
					\item Similarly. 
				\end{itemize}
			\end{proof}
			
			
				
			\begin{proposition}
				Let: $\mathcal{A}_0$ be a semialgebra. 
				
				Then, every additive set function $\mu$ on $\mathcal{A}_0$ uniquely extends to an additive set function on the algebra $\sigma(\mathcal{A}_0)$. 
				
				This extension is countably additive provided that $\mu$ is countably additive on $\mathcal{A}_0$. 
				
				The same is true for semiring and ring. 
			\end{proposition}
			\begin{proof}
				\begin{itemize}
					\item The extension: it is clear, that every element in $\mathcal{A}$ could be represented as the union of disjointed elements of $\mathcal{A}_0$. Now, set
					\begin{gather*}
						\mu(A):=\sum_{i=1}^n \mu(A_i)\qquad A_i\in \mathcal{A}_0
					\end{gather*}
					
					Then, it is obviously additive. And such summation is idependent of choice of partion: If $B=\sqcup B_i$, then consider the equality
					\begin{gather*}
						\mu(A_i)=\sum \mu(A_i\cap B_j)\qquad \mu(B_j)=\sum \mu(A_i\cap B_j)
					\end{gather*}
					\item Now, the countable additivity: Suppose 
					\begin{gather*}
						A=\cup_{n=1}^\infty A_n\qquad A_n\cap A_k=\emptyset,n\not=k\qquad A_n\in \mathcal{A}
					\end{gather*}
					
					Then, for a family $A=\cup_{j=1}^N B_j,A_n=\cup_{i=1}^{N_n} B_{n,i}$, where $B_j,B_{n,j}\in \mathcal{A}_0$, set
					\begin{gather*}
						C_{n,i,j}=B_{n,i}\cap B_j\in \mathcal{A}_0
					\end{gather*}
					
					Then, $\{c_{n,i,j}\}$ are pairwise disjoint, and then
					\begin{gather*}
						B_j=\cup_{n=1}^\infty \cup_{i=1}^{N_n} \mu(C_{n,i,j}),B_{n,j}=\cup_{i=1}^{N_n} C_{n,i,j} \Rightarrow \mu(B_j)=\sum_{n=1}^\infty \sum_{i=1}^{N_n} \mu(C_{n,i,j}),\mu(B_{n,j})=\sum_{i=1}^{N_n} \mu(C_{n,i,j})
					\end{gather*}
					
					Finally,
					\begin{gather*}
						\mu(A)=\sum_{j=1}^N \sum_{n=1}^\infty \sum_{i=1}^{N_n} \mu(C_{n,i,j})=\sum_{n=1}^{\infty}\sum_{j=1}^N \sum_{i=1}^{N_n} \mu(C_{n,i,j})
					\end{gather*}
					\item Similarly for ring. 
				\end{itemize}
			\end{proof}
			
			
			\subsubsection{Probability}
			
			Especially: 
			
			\begin{Def}
				\begin{itemize}
					\item A \textbf{countably additive measure} $\mu$ on a $\sigma$-algebra is called a \textbf{probability measure} if
					\begin{gather*}
						\mu\geq 0\qquad \mu(X)=1
					\end{gather*}
					
					\item Such $(X,\mathcal{A},\mu)$ is called a \textbf{probability space}.
				\end{itemize}
			\end{Def}
			\subsubsection{Compact Classes, and Countable Additivity}
			\begin{Def}
				A family $K$ of subsets of a set $X$ is called \textbf{compact} if, for any sequences $K_n$ of its elements, 
				\begin{gather*}
					\cap_{n=1}^\infty K_n\Rightarrow \cap_{n=1}^N K_n
				\end{gather*}
			\end{Def}
			
			\begin{proposition}
				Let $\mu$ be a nonnegative additive set function on an algebra $\mathcal{A}$. 
				
				Suppose that: There exists a compact class $\mathcal{K}$ approximating $\mu$ in following sense: For every $A\in \mathcal{A}$, there exists 
				\begin{gather*}
					\exists K_\epsilon\in \mathcal{K},A_\epsilon\subseteq K_\epsilon\subseteq A\qquad \mu(A-A_\epsilon)<\epsilon
				\end{gather*}
				
				Then, $\mu$ is countably additive. In particular, this is true if the compact class $\mathcal{K}$ is contained in $\mathcal{A}$ and for any $A\in \mathcal{A}$, one has the equality
				\begin{gather*}
					\mu(A)=\sup_{K\subseteq A,K\in \mathcal{K}} \mu(K)
				\end{gather*}
			\end{proposition}
			\begin{proof}
				Suppose that the sets $A_n\in \mathcal{A}$ are decreasing, and 
				\begin{gather*}
					\cap_n A_n=\emptyset
				\end{gather*}
				
				Then, $\mu(A_n)\rightarrow 0$: Suppose that there exists a family of $K_n,B_n$ such that $B_n\subseteq K_n\subseteq A_n,\mu(A_n-B_n)<2^{-n}\epsilon$.
				
				By compactness, there exists some $N$, $\cap_{n=1}^N K_n=\emptyset\Rightarrow \cap_{n=1}^N B_n=\emptyset$. THen, one has
				\begin{gather*}
					A_N=\cap_{n=1}^N A_n\subseteq \cup_{n=1}^N (A_n-B_n)\\
					\Rightarrow \mu(A_n)\leq \sum_{n=1}^N \epsilon 2^{-n}\epsilon\leq \epsilon
				\end{gather*}
			\end{proof}
			
			\subsubsection{Approxiamating Class}
				\begin{Def}
					Let $m$ be a nonnegative function on a class $\mathscr{E}\subseteq P(X)$, and $\mathcal{P}$ be a class of subsets of $X$. 
					
					We say that $P$ is an approximating class for $m$, if for every $E\in \mathscr{E}$, every $\epsilon>0$, there exists 
					\begin{gather*}
						P_\epsilon\in \mathcal{P},E_\epsilon\in \mathcal{E};E_\epsilon\subseteq P_\epsilon\subseteq E,|m(E)-m(E_\epsilon)|<\epsilon
					\end{gather*}
				\end{Def}
				
			\subsubsection{Approxiamating Class on Borel Set}
			
				\begin{lemma}
					Consider the class
					\begin{gather*}
						\mathcal{A}=\{A\in \mathcal{B}(\RR^n)|\forall \epsilon>0, F\mbox{ closed},U\mbox{ open},F\subseteq A\subseteq U,\mu(U-F)<\epsilon\}
					\end{gather*}
					
					Then, $\mathcal{A}=B(\RR^n)$
				\end{lemma}
				\begin{proof}
					This is a Borel $\sigma$-algebra: 
					
					\begin{itemize}
						\item Firstly, the closed subsets belongs to $\mathcal{A}$: Easy to check, just take the union of disks $A^\delta=\{x\in \RR^n|dist(x,A)<\delta\}$, notice that
						\begin{gather*}
							\lim_{k\rightarrow\infty} (A^\frac{1}{k}-A)=0
						\end{gather*}
						
						Thus, we could choose the approxiamating to $A$.
						
						\item Secondly, $\mathcal{A}$ is stable under compeletion:
						\begin{gather*}
							F\subseteq A\subseteq U\Rightarrow U^c\subseteq A^c\subseteq F^c
						\end{gather*}
						
						\item Thirdly: $\mathcal{A}$ is stable under countable union:
						\begin{gather*}
							\mu(U_j-F_j)<\frac{\epsilon}{2^{j+1}}\Rightarrow \mu(U-F)\leq \sum \mu(U_j-F_j)<\frac{\epsilon}{2}
						\end{gather*} 
						
						But $F=\cup F_j$ may be not closed. Take a smaller $Z_k=\cup_{1}^k F_j$. Then, suppose $k$ is large enough, so $\mu(F-Z_k)<\frac{\epsilon}{2}$,
						\begin{gather*}
							\mu(U-Z_k)\leq \mu(U-F)+\mu(F-Z_k)<\epsilon
						\end{gather*}
					\end{itemize}
					
					Proved. 
				\end{proof}
			
				\begin{proposition}
					(Regularity of Borel Measure)Let $\mu$ be a nonnegative contably additive measure on the Borel $\sigma$-algebra $B(\RR^n)$ in the space $\RR^n$. 
					
					Then, for any Borel set $B\subseteq \RR^n$, and $\epsilon>0$, there exist an open subset $U_\epsilon$ and a compact set $K_\epsilon$ such that
					\begin{gather*}
						K_\epsilon\subseteq B\subseteq U_\epsilon\qquad \mu(U_\epsilon-K_\epsilon)<\epsilon
					\end{gather*}
				\end{proposition}
				\begin{proof}
					We've known, that for Borel set, we have approxiamation
					\begin{gather*}
					F\subseteq B\subseteq U\qquad \mu(U-F)<\frac{\epsilon}{2}
					\end{gather*}
					
					where $F$ is closed. Then select $K_R=F\cap \bar{B}_R$, which is bounded and closed, so compact. 
					
					Notice that, $\lim_{R\rightarrow \infty} K_R=F$. Then, 
					\begin{gather*}
						\mu(F-K_R)\rightarrow 0
					\end{gather*}
					
					Choose such $R$ with $\mu(F-K_R)<\frac{\epsilon}{2}$, we obtrain 
					\begin{gather*}
						K_R\subseteq B\subseteq U
					\end{gather*}
				\end{proof}
				
				\begin{corollary}
					$B(\RR^n)$ is measurable. 
				\end{corollary}
			\subsection{Infinite Measure,$\sigma$-Finite Measures}	
				Now, we will extend the value of measure to $[0,+\infty]$:
				
				\begin{Def}
					Let 
					\begin{itemize}
						\item $\mathcal{A}$ be a $\sigma$-algebra in a space $X$, and 
						\item  let $\mu$ be a set function on $\mathcal{A}$, with values in $[0,+\infty]$, and satisfies
						\begin{itemize}
							\item $\mu(\emptyset)=0$;
							\item $\mu$ is countably additive
							\begin{gather*}
								\mu(\cup A_j)=\sum \mu(A_j)
							\end{gather*}
						\end{itemize}
						
						Then, $\mu$ is called a measure with values in $[0,\infty]$;
						
						Especially, $\mu$ is called a \textbf{$\sigma$-finite} measure, if $X=\cup_{n=1}^\infty X_n,X_n\in \mathcal{A},\mu(X_n )<\infty$;
					\end{itemize}
				\end{Def}
				
				\begin{lemma}
					Let $\mathcal{R}$ be a ring of subsets of space $X$. 
					
					Let $\mu$ be a contably additive function on $\mathcal{R}$ with value $[0,+\infty]$ such that there exist sets $X_n\in \mathcal{R}$ with $X=\cup X_n,\mu(X_n)<\infty$. 
					
					Denote: $\mu_n$ the \underline{Lesbegue extension} of measure $\mu$, regarded on the sets 
					\begin{gather*}
						S_n=\cup_{j=1}^n X_j
					\end{gather*} 
					
					equipped with the algebra of sets consisting of the intersection of elements in $\mathcal{R}$ with $S_n$. Let $\mathcal{L}_{\mu_n}$ denote the class of all $\mu_n$-measurable sets. Let
					\begin{gather*}
						\mathcal{A}:=\{A\subseteq X|A\cap S_n\in \mathcal{L}_{\mu_n},\forall n\in \NN,\}
					\end{gather*}
					
					Then, $\mathcal{A}$ is a ring with respect to countable intersection, and $\bar{\mu}$ is a $\sigma$-additive measure, whose restriction to those $S_n$ coincide with $\mu$. 
				\end{lemma}
				\begin{proof}
					Just consider 
					\begin{gather*}
						\mu_n(A\cap S_n )=\sum \mu_n(A_i\cap S_n)
					\end{gather*}
				\end{proof}
				
				\begin{proposition}
					Let $\mu$ be a countably additive additive measure on an algebra $\mathcal{A}$, with values in $[0,\infty]$, then $A_\mu$ is a $\sigma$-algebra, $A\subseteq A_\mu$, and the function $\mu^*$ is a countably additive measure on $\mathcal{A}_\mu$ with values in $[0,+\infty]$ and coincides with $\mu$ on $\mathcal{A}$. 
				\end{proposition}
				\begin{proof}
					内容...
				\end{proof}
		\section{Outer Measure}
			\subsection{Outer Measure and the Lebesgue Extension Measures}
				\begin{Def}
					Let: $\mu$ be a nonnegative set function $\mu$ that is defined a certain class $\mathcal{A}$ of subsets in a space $X$, and contains $X$ itself, then the formula
					\begin{gather*}
						\mu^*(A)=\inf\{\sum_{n=1}^\infty \mu(A_n)|A_n\in \mathcal{A},A\subseteq \cup_{n=1}^\infty A_n\},A\subseteq X
					\end{gather*}
					
					Then, $\mu^*:P(X)\rightarrow [0,+\infty]$. 
				\end{Def}
				
				
				\begin{Def}
					Suppose that $\mu$ is nonnegative set function on domain $\mathcal{A}\subseteq 2^X$. 
					
					A set $A$ is called \textbf{$\mu$-measurable} if for any $\epsilon>0$, there exists $A_\epsilon\in \mathcal{A}$, such that
					\begin{gather*}
						\mu^*(A\Delta A_\epsilon)<\epsilon
					\end{gather*}
					
					Take the notion of $\mu$-measurable subsets $A_\mu$. 
				\end{Def}
				
				\begin{lemma}
					Let $\mu$ be a nonnegative set function on a class $\mathcal{A}$, then the function $\mu^*$ is contably subadditive. i.e.
					\begin{gather*}
						\mu^*(\cup_{n=1}^\infty A_n)\leq \sum_{n=1}^\infty \mu^*(A_n)
					\end{gather*}
					
					for any sets $A_n$. 
				\end{lemma}
				\begin{proof}
					内容...
				\end{proof}
				
				\begin{lemma}
					In the situation of previous lemma, for any sets $A$ and $B$, such that $\mu^*(B)<\infty$, one has the equality
					\begin{gather*}
						|\mu^*(A)-\mu^*(B)|\leq \mu^*(A\Delta B)
					\end{gather*}
				\end{lemma}
				\begin{proof}
					内容...
				\end{proof}
				
				\begin{proposition}
					Let $\mu$ be a nonnegative countably additive set function on an algebra $\mathcal{A}$, then:
					\begin{itemize}
						\item One has $\mathcal{A}\subseteq \mathcal{A}_\mu$, and the outer measure $\mu^*$ coincides with $\mu$ on $\mathcal{A}$;
						\item The collection $\mathcal{A}_\mu$ of all $\mu$-measurable sets is a $\sigma$-algebra, nad the striction of $\mu^*$ to $\mathcal{A}_\mu$ is countably additive;
						\item The function $\mu^*$ is a unique nonnegative countably additive extension of $\mu$ to the $\sigma$-algebra $\sigma(\mathcal{A})$ generated by $\mathcal{A}$, and a unique nonnegative countably additive extension of $\mu$ to $\mathcal{A}_\mu$. 
					\end{itemize}
				\end{proposition}
				\begin{proof}
					内容...
				\end{proof}
				
				\begin{corollary}
					Let $\mu$ be a nonnegative countably additive set function on an algebra $\mathcal{A}$.
					
					A set $A$ is $\mu$-measurable precisely when there exist two sets $A',A''\in \sigma(\mathcal{A})$ such that
					\begin{gather*}
						A'\subseteq A\subseteq A''\qquad \mu^*(A''-A')=0
					\end{gather*}
					
					Moreover, one can take for $A'$ a set of the form $\cup_{n=1}^\infty \cap_{k=1}^\infty A_{n,k},A_{n,k}\in\mathcal{A}$, and for $A''$ a set of the form $\cap_{n=1}^\infty\cup_{k=1}^\infty B_{n,k},B_{n,k}\in \mathcal{A}$. 
				\end{corollary}
				\begin{proof}
					内容...
				\end{proof}
				
				\begin{corollary}
					内容...
				\end{corollary}
				\begin{proof}
					内容...
				\end{proof}
				
				\subsubsection{Collection $\mathcal{A}_\mu$}
					\begin{proposition}
						Let $\mu$ be a nonnegative countably additive measure on an algebra $\mathcal{A}$. Then, TFAE:
						\begin{itemize}
							\item $A\in \mathcal{A}_\mu$;
							\item $\mu^*(A)+\mu^*(X-A)=\mu^*(X)$;
							\item $\mu^*(E\cap A)+\mu^*(E-A)=\mu^*(E),E\subseteq X$. 
						\end{itemize}
					\end{proposition}
					\begin{proof}
						\begin{itemize}
							\item 
							
							\item 
							
							\item 
						\end{itemize}
					\end{proof}
				\subsubsection{Countable Additivity}
					\begin{proposition}
						Let $\mu$ be a nonnegative measure on a $\sigma$-algebra $\mathcal{A}$, suppose that sets $\mathcal{A}_n$ are such that $A_n\subseteq A_{n+1}$ for all $n\in \NN$, then, one has
						\begin{gather*}
							\mu^*(\cup A_n)=\lim \mu^*(A_n)
						\end{gather*}
					\end{proposition}
					\begin{proof}
						内容...
					\end{proof}
				\subsubsection{Compeleteness}
					\begin{Def}
						内容...
					\end{Def}
					
				
				
			\subsection{Carathedory Outer Measure}
				\begin{Def}
					A set function $m:P(X)\rightarrow [0,+\infty]$ defined on the class of all subsets of a set, and taking values in $[0,+\infty]$ is called \textbf{outer measure} on $X$, if:
					\begin{itemize}
						\item $m(\emptyset)=0$;
						\item $m(A)\leq m[B],A\subseteq B$;
						\item $m(\cup_{n=1}^\infty A_n)\leq \sum_{n=1}^\infty m(A_n)$;
					\end{itemize}
				\end{Def}
				
				Especially,
				\begin{Def}
					Let $m$ be a set function with values in $[0,+\infty]$ defined on the class of all subsets of a space $X$, such that $m(\emptyset)=0$. 
					
					A set $A\subseteq X$ is called \textbf{Caratheodory measurable} with respect to $m$, if for every $E\subseteq X$, one has the equality
					\begin{gather*}
						m(E\cap A)+m(E-A )=m(E)
					\end{gather*}
					
					The class of all Caratheodory $m$-measurable sets is denoted by $\mathfrak{M}_m$. 
				\end{Def}
				
				\begin{proposition}
					Let $m$ be a set function with values in $[0,+\infty]$ on the class of all sets in a space $X$, such that $m(\emptyset)=0$, then
					\begin{itemize}
						\item 
						\item 
						\item 
					\end{itemize}
				\end{proposition}
				\begin{proof}
					内容...
				\end{proof}
					\section{Lesbegue Measure}
		\subsection{Lesbegue Measure}
		\begin{Def}
			\begin{itemize}
				\item 
				Consider: The \textbf{cube} $I$ in $\RR^n$ of form $[a,b]^n$
				\item Then, the algebra
				\begin{gather*}
					\mathcal{A}_0=\{\cup_j=1^m I_j|m\in \NN,I_j\mbox{ is cube}\}
				\end{gather*}
				\item Volumes:
				\begin{gather*}
					\lambda^0_n(\prod [a_i,b_i)):=\prod (b_i-a_i)
				\end{gather*}
				
				Then, by finite additivity, we could extend it to the whole $\mathcal{A}_0$.
				
				Especially: Suppose $\sqcup_{j=1}^\infty A_j\in \mathcal{A}_0$, then
				\begin{gather*}
					内容...
				\end{gather*}
				
				Thus, this is a \underline{premeasure}
				
				\item According the construction by Caratheodory, we got an outer measure $\lambda^*_n$. This is the \textbf{Lesbegue $\sigma$-Algebra} on $I$.
				
				\item Finally, define:
				\begin{gather*}
					内容...
				\end{gather*}
				
				By the increasing cubes, 
			\end{itemize}
		\end{Def}
		\subsection{Lesbegue-Stieltjes Measures}
		\begin{Def}
			(Distribution Function)
			\begin{itemize}
				\item Suppose $\mu$ be a nonnegative Borel measure on $\RR$. Then, the function
				\begin{gather*}
					t\mapsto F(t)=\mu((-\infty,t))
				\end{gather*}
				
				is bounded, nondecreasing
				
				\item Such function is often called the \textbf{distribution function}.
			\end{itemize}
		\end{Def}
		
		\begin{theorem}
			Let $F$ be a bounded, nondecreasing, left continuous function, with
			\begin{gather*}
				\lim_{t\rightarrow-\infty} F(t)=0
			\end{gather*}
			
			Then, there exists a unique nonnegative Borel measure on $\RR$ such that
			\begin{gather*}
				F(t)=\mu((-\infty,t))\qquad \forall t\in \RR
			\end{gather*}
		\end{theorem}
		\begin{proof}
			内容...
		\end{proof}
		\section{Monotone and $\sigma$-additive classes of sets}
		
		\section{Souslin Sets, and $A$-Operations}
			\begin{Def}
				内容...
			\end{Def}
		
	\chapter{The Lesbegue Integral}
		\section{Measurable Functions}
			\begin{Def}
				Let $(X,\mathcal{A})$ be a measurable space.
				
				Then, a function $f:X\rightarrow \RR$ is called \textbf{measurable} if
				\begin{gather*}
					\{x|f(x)<c\}\subseteq X\in \mathcal{A}\qquad \forall c\in \RR
				\end{gather*}
			\end{Def}
			
			\begin{proposition}
				TFAE:
				\begin{itemize}
					\item 
					A function $f$ is \textbf{measurable} with respect to a $\sigma$-algebra $\mathcal{A}$ 
					\item $f^{-1}(B)\in \mathcal{A},\forall B\in \mathscr{B}(\RR)$. 
				\end{itemize}
			\end{proposition}
			\begin{proof}
				内容...
			\end{proof}
			
			\begin{Def}
				Define: The \textbf{indicator}
				\begin{gather*}
					1_A(x):=\begin{cases}
						1&x\in A_i\\
						0&Elsewise
					\end{cases}
				\end{gather*}
				
				It is easy to verify, $1_A$ is measurable. 
			\end{Def}
			
			Now, take the following notion
			\begin{gather*}
				f^+(x):=\max(f(x),0)\qquad f^-(x):=\max(-f(x),0)
			\end{gather*}
			
			Then,
			\begin{proposition}
				TFAE: For function $f:\mathcal{A}\rightarrow \RR$
				\begin{itemize}
					\item $f$ is measurable;
					\item $f^+,f^-$ are measurale;
				\end{itemize}
			\end{proposition}
			\begin{proof}
				内容...
			\end{proof}
			\begin{proposition}
				TFAE： For function $f:(X,\mathcal{A})\rightarrow (\RR,B(\RR))$,
				\begin{itemize}
					\item $f$ is measurable;
					\item $f^{-1}(B)\in \mathcal{A},B\in \mathcal{B}(\RR)$;
				\end{itemize}
			\end{proposition}
			\begin{proof}
				内容...
			\end{proof}
			\subsection{Mesurable Function beetween Measurable Spaces}
			
			\begin{Def}
				Let $(X_1,\mathcal{A}_1),(X_2,\mathcal{A}_2)$ be two spaces with $\sigma$-algebras. 
				
				A mapping $f:X_1\rightarrow X_2$ is called \textbf{measurable} with respect to the pari $(\mathcal{A}_1,\mathcal{A}_2)$ if
				\begin{gather*}
					B\in \mathcal{A}_2\Rightarrow f^{-1}(B)\in \mathcal{A}_1
				\end{gather*}
			\end{Def}
			
			\begin{proposition}
				Suppose that the function $f,g,f_n:X\rightarrow \RR$, where $n\in \NN$, are measurable with respect to a $\sigma$-algebra $\mathcal{A}$, then
				\begin{itemize}
					\item The function $\varphi\circ f$ is measurable with respect to $\mathcal{A}$ for Borel function
					\begin{gather*}
						\varphi:\RR\rightarrow \RR
					\end{gather*}
					
					in particular, this is true if $\varphi$ is continuous. 
					\item The function
					\begin{gather*}
						\alpha f+\beta g
					\end{gather*}
					
					is measurable with respect to $\mathcal{A}$, for every $\alpha,\beta\in \RR$;
					\item The function $fg$ is measurable with respect to $\mathcal{A}$;
					\item If $g(x)\not=0$, the function $\frac{f}{g}$ is measurable with respect to $\mathcal{A}$;
					\item If there exists a finite limit $f_0(x)=\lim_{n\rightarrow \infty} f_n(x)$ for all $x$, then the function $f_0$ is measurable with respect to $\mathcal{A}$;
					\item If the functions $\sup_n f_n(x)$ and $\inf_n f_n(x)$ are finite for all $x$, then they are measurable with respect to $\mathcal{A}$. 
				\end{itemize}
			\end{proposition}
			\begin{proof}
				\begin{itemize}
					\item 
					\item 
					\item 
					\item 
					\item 
					\item 
				\end{itemize}
			\end{proof}
			
			\begin{proposition}
				内容...
			\end{proposition}
			\begin{proof}
				内容...
			\end{proof}
			
			\subsubsection{$\mu$-Measurable}
			
			\begin{Def}
				Let $(X,\mathcal{A},\mu)$ be a measure space. 
				
				A real function $f$ on $X$ is called \textbf{$\mu$-measurable} if it is measurable with respect to the $\sigma$-algebra $\mathcal{A}_\mu$ of all $\mu$-measurable sets.
				
				In addition, we agree to call $\mu$-measurable aloso any function $f$ that is defined and $\mathcal{A}_\mu$-measuralbe on $X-Z$, where $Z$ is a set of $\mu$-measure zero. 
				
				The set of all $\mu$-measurable function is denoted by $\mathcal{L}^0(\mu)$. 
			\end{Def}
			
			\begin{proposition}
				Let $\mu$ be a nonnegative measure on a $\sigma$-algebra $\mathcal{A}$. 
				
				Then, for every $\mu$-measurable function $f$, one can find a set $Y\in \mathcal{A}$ and a function $g$ measurable with respect to $\mathcal{A}$ such that
				\begin{gather*}
					f(x)=g(x)\qquad \forall x\in Y,\mu(X-Y)=0
				\end{gather*}
			\end{proposition}
			\begin{proof}
				内容...
			\end{proof}
			
			
			\subsection{Convergence, Almost Everywhere }
			
			\begin{Def}
				\begin{itemize}
					\item Some property for points in $X$ is fullfilled \textbf{$\mu
					$-almost everywhere} on $X$ if the set $Z$ of all points in $X$ that do not have this property 
					\begin{itemize}
						\item belongs to $A_\mu$, and
						\item has measure zero with respect to $\mu$. 
					\end{itemize}
					
					Denoted by 'a.e.'
					\item $f=g$ a.e. states, that
					\begin{gather*}
						\{x|f(x)=g(x)\}\in \mathcal{A},\mu\{x|f(x)=g(x)\}=0
					\end{gather*}
				\end{itemize}
			\end{Def}
			
			
			\begin{theorem}
				(Egoroff) Let $(X,\mathcal{A},\mu)$ be a space with a finite nonnegative measure $\mu$, and let $\mu$-measurable function $f_n$ be such that $\mu$-almost everywhere, there is a finite limit
				\begin{gather*}
					f(x)=\lim_{n\rightarrow\infty} f_n(x)
				\end{gather*}
				
				Then for every $\epsilon>0$, there exists a set $X_\epsilon\in \mathcal{A}$, such that
				\begin{itemize}
					\item $\mu(X-X_\epsilon)<\epsilon$;
					\item the function $f_n\stackrel{uniformly}{\rightarrow} X_\epsilon$ on $X_\epsilon$. 
				\end{itemize}
			\end{theorem}
			\begin{proof}
				Define:
				\begin{gather*}
					X^m_n:=\cap_{i\geq n} \{x||f_i(x)-f(x)|<\frac{1}{m}\}\in \mathcal{A}_\mu
				\end{gather*}
				
				Then, $X^m_n\subseteq X_{n+1}^m$. And $\cup_{n=1}^\infty X_n^m=X$: This is because that, for every $y\in X$, and fixed $m\in \ZZ^+$, because
				\begin{gather*}
					\lim_{n\rightarrow\infty} f_n(y)=f(y)
				\end{gather*}
				
				there exists large enough $n$ such that $y\in X^m_n$. 
				
				Thus, for each $\epsilon>0$, there exists a number $k(m)$ such that
				\begin{gather*}
					\mu(X-X_{k(m)}^m)<\epsilon 2^{-m}
				\end{gather*}
				
				Now, let $X_\epsilon=\cap_{m=1}^\infty X_{k(m)}^m$, we obtain
				\begin{gather*}
					\mu(X-X_\epsilon)\leq \sum \mu(X-X_{k(m)}^m)\leq \epsilon\sum_m 2^{-m}=\epsilon
				\end{gather*}
			\end{proof}
			
			\begin{Def}
				Suppose we are given a measure space $(X,\mathcal{A},\mu)$ with a finite measure $\mu$, and a sequence of $\mu$-measurable functions:
				\begin{itemize}
					\item The sequence $\{f_n\}$ is called \textbf{fundamental} or \textbf{Cauchy} if, for every $c>0$, one has
					\begin{gather*}
						\lim_{N\rightarrow\infty} \sup_{n,k\geq N} \mu(x||f_n(x)-f_k(x)|\geq c)=0
					\end{gather*}
					\item The sequence $\{f_n\}$ is said to \textbf{converge} in measure to a $\mu$-measurable function $f$ if, for every $c>0$, one has
					\begin{gather*}
						\lim_{n\rightarrow\infty} \mu(x||f(x)-f_n(x)|\geq c)=0
					\end{gather*}
				\end{itemize}
			\end{Def}
			
			Notice, $	\lim_{N\rightarrow\infty} \sup_{n,k\geq N} \mu(x||f_n(x)-f_k(x)|\geq c)$ is well-defined, because it is a descending sequence. 
			
		\begin{proposition}
			Let $(X,\mathcal{A},\mu)$ be a space with a \underline{finite measure}($\mu(X)<\infty$)
			
			\begin{itemize}
				\item If a sequence of $\mu$-measurable functions $f_n$ converges to $f$ in measure $\mu$, then there exists its subsequence $f_{n_k}\rightarrow f$ almost everywhere. 
				\item If a sequence of $\mu$-measurable functions $f_n$ is fundamental in measure $\mu$, then it cconverges in measure $\mu$ to some measurable function $f$. 
			\end{itemize}
		\end{proposition}
		\begin{proof}
			\begin{itemize}
				\item 
				\item 
			\end{itemize}
		\end{proof}
		
		\begin{corollary}
			Let $\mu$ be a finite measure and let two sequences of $\mu$-measurable function $f_n,g_n$ converge in measure $\mu$ to functions $f$ and $g$.
			
			Suppose that $\Psi$ is a continuous function on $Y\subseteq \RR^2$, such that
			\begin{gather*}
				(f(x),g(x))\in Y\qquad (f_n(x),g_n(x))\in Y
			\end{gather*}
			
			Then, $\Psi(f_n,g_n)\rightarrow \Psi(f,g)$.
		\end{corollary}
		\begin{proof}
			内容...
		\end{proof}
		
		\begin{corollary} For such $f_n,g_n,f,g$,
			\begin{gather*}
				f_ng_n\rightarrow fg\qquad \alpha f_n+\beta g_n\rightarrow \alpha f+\beta g
			\end{gather*}
		\end{corollary}
		\begin{proof}
			内容...
		\end{proof}
		
		\subsubsection{Approxiamation Extension from Compact}
			\begin{lemma}
				内容...
			\end{lemma}
			\begin{proof}
				内容...
			\end{proof}
			
			\begin{proposition}
				内容...
			\end{proposition}
			\begin{proof}
				内容...
			\end{proof}
			
			\begin{theorem}
				(Lusin's Theorem) A function $f$ on an interval $I$, with Lesbegue measure is measurable precisely when for each $\epsilon>0$, there exists a continuous function 
				\begin{gather*}
					f_\epsilon:\RR\rightarrow \RR
				\end{gather*}
				
				and a compact set $K_\epsilon$ such that
				\begin{gather*}
					\mu(I-K_\epsilon)<\epsilon\qquad f=f_\epsilon\mbox{ on }K_\epsilon
				\end{gather*}
			\end{theorem}
			\begin{proof}
				内容...
			\end{proof}
		\section{Integrations}
			\subsection{Simple Functions}
		
			\begin{Def}
				\begin{itemize}
					\item 
					Let $(X,\mathcal{A})$ be a measurable space.
					
					A function
					\begin{gather*}
						f:X\rightarrow \RR
					\end{gather*}
					
					is called \textbf{simple function}, if $f$ is measurable, and $f(X)$ is a finite subset. 
				\end{itemize}
			\end{Def}
			
			\begin{proposition}
				TFAE:
				\begin{itemize}
					\item $f$ is a simple function;
					\item $f=\sum a_i 1_{A_i},\cup A_i\subseteq A$;
					\item $f=\sum a_i 1_{A_i}.\sqcup A_i\subseteq A$;
				\end{itemize}
			\end{proposition}
			\begin{proof}
				内容...
			\end{proof}
			
			\begin{Def}
				内容...
			\end{Def}
			
			\begin{lemma}
				Let $\{f_n\},\{g_n\}$ be tow sequences of simple functions, that are:
				\begin{itemize}
					\item mean fundamental($L^1$ Cauchy)
					\item Almost everywhere to one and the same function $f$.
				\end{itemize}
				
				Then, $f_n$ and $g_n$ converge to the same value. 
			\end{lemma}
			\begin{proof}
				内容...
			\end{proof}
			
			\begin{lemma}
				内容...
			\end{lemma}
			\begin{proof}
				内容...
			\end{proof}
			
			\begin{Def}
				内容...
			\end{Def}
			
			\subsection{Integral of Simple Functions}
			
			\begin{Def}
				Let $(X,\mathcal{A},\mu)$ be a space with finite nonnegative measure. 
				
				For any simple function $f$ on $X$, that assmues finitely many values $c_i$ on disjoint sets $A_i$, the \textbf{Lesbegue integral} of $f$ with respect to $\mu$ is defined by
				\begin{gather*}
					\int_X f(x)\mu (dx):=\sum_i c_i\mu(A_i)
				\end{gather*}
			\end{Def}
			
			\subsubsection{Fundamental/Mean Fundamental/Fundamental in $L^1(\mu)$}
			\begin{Def}
				内容...
			\end{Def}
			
			\begin{lemma}
				The Lesbegue integral on simple functions enjoys the following properties:
				\begin{itemize}
					\item 
					\item 
					\item 
				\end{itemize}
			\end{lemma}
			\begin{proof}
				内容...
			\end{proof}
			
			\begin{corollary}
				内容...
			\end{corollary}
			\begin{proof}
				内容...
			\end{proof}
			
			\begin{proposition}
				内容...
			\end{proposition}
			\begin{proof}
				内容...
			\end{proof}
			
			\subsection{Lesbegue Integration}
			
			\begin{Def}
				内容...
			\end{Def}
			
			
			\subsubsection{Basic Properties}
			
			\subsubsection{Chebyshev's Theorem}
			
			\begin{theorem}
				(Chebyshev/Markov Inequality) 
				\begin{gather*}
					\mu(\{x||f(x)|\geq R\})\leq \frac{1}{R}\int_X |f(x)|\mu(dx)\qquad R>0
				\end{gather*}
			\end{theorem}
			\begin{proof}
				内容...
			\end{proof}
			
			\begin{corollary} 
				\begin{gather*}
					\mu(\{|f|\geq \lambda\})\leq \frac{\Vert f\Vert^p_{L^p}}{\lambda^p}
				\end{gather*}
			\end{corollary}
			\begin{proof}
				内容...
			\end{proof}
			\subsubsection{Covergence Theorem:Lesbegue/Levi/Fatou}
				
				\begin{theorem}
					(Lesbegue Dominated Covergence Theorem) Suppose that $\mu$-integrable functions $f_n$ converge almost everywhere to a function $f$.
					
					If there exists a $\mu$-integrable function $\Phi$, such that
					\begin{gather*}
						|f_n(x)|\leq \Phi(x),\forall n
					\end{gather*}
					
					Then, 
					\begin{itemize}
						\item the function $f$ is integrable, and 
						\begin{gather*}
							\int_X f(x)\mu(dx)=\lim_{n\rightarrow\infty} f_n(x)\mu(dx)
						\end{gather*}
						\item In addition
						\begin{gather*}
							\lim_{n\rightarrow\infty} \int_X|f(x)-f_n(x)|\mu(dx)=0
						\end{gather*}
					\end{itemize}
				\end{theorem}
				\begin{proof}
					内容...
				\end{proof}
			
			\subsection{Criteria of Integrability}
				\begin{theorem}
					内容...
				\end{theorem}
				\begin{proof}
					内容...
				\end{proof}
		\section{Space $L^1$}
			
			\subsection{Completeness}
			\begin{Def}
				\begin{itemize}
					\item A sequence of function $f_n$ that are integrable with respect to a measure $\mu$, is calledd \textbf{fundamental}
					\item 
				\end{itemize}
			\end{Def}
		\section{Riemann Integral}
	\chapter{Operations on Measures and Functions}
		\section{Decompositikon on Measures, and Functions}
	\chapter{Integral and Derivative}
		\section{Differentiability of Functions on the Real Line}
		
	\chapter{Borel,Baire and Souslin Sets}
	\chapter{Measures on Topological Spaces}
		\section{Meaasures on Manifolds}
	\chapter{Weak Convergence of Measure}
	\chapter{Transformations of Measures, and Isomorphisms}
	\chapter{Conditional Measures and Conditional Expectation}
		
\end{document}